\chapter{估计、检验、渐近、回归与多重检验}
\label{ch:statistical-inference}

观察到十次试验中七次成功，怎样估计下一次成功的概率？用一条直线解释带噪观测，为什么会得到最小二乘？本章先从分布假设构造估计量，再研究它的抽样误差、置信区间和检验，最后讨论回归与多重搜索。

同一个估计值在不同抽样模型下可能具有不同的精度。因此，构造估计量与评价估计量是相连的两个问题。

\section{估计量与抽样误差}

设 $X_1,\ldots,X_n\sim P_\theta$，估计量
$\hat\theta_n=T_n(X_1,\ldots,X_n)$ 是随机变量。它的偏差与均方误差定义为
\begin{align*}
 \operatorname{Bias}_\theta(\hat\theta_n)
 &=\mathbb E_\theta\hat\theta_n-\theta,\\
 \operatorname{MSE}_\theta(\hat\theta_n)
 &=\mathbb E_\theta(\hat\theta_n-\theta)^2
 =\operatorname{Var}_\theta(\hat\theta_n)
  +\operatorname{Bias}_\theta(\hat\theta_n)^2.
\end{align*}
不需要强调参数与样本量时，也简写为
$\operatorname{Bias}(\hat\theta)$。
最后一个等式不是定义。令
$b=\mathbb E\hat\theta_n-\theta$，把误差写成
$(\hat\theta_n-\mathbb E\hat\theta_n)+b$；平方取期望后，交叉项因中心化为零，便得到分解。

\MFQLead{有偏估计也可能更好}
若无偏估计量 $\bar X$ 估计 $\theta=1$，且
$\operatorname{Var}(\bar X)=0.25$，考虑收缩估计
$\tilde\theta=0.8\bar X$。它的偏差为 $-0.2$、方差为
$0.8^2\times0.25=0.16$，所以 MSE 是
$0.16+(-0.2)^2=0.20$，小于无偏估计的 $0.25$。
无偏不是“最佳”的同义词；比较估计量必须先声明损失函数与参数范围。

反过来，无偏也不保证一致。若 $X_i\sim\mathcal N(\theta,1)$，取
$T_n=X_1$，则对每个 $n$ 都有 $\mathbb ET_n=\theta$，但
$\operatorname{Var}(T_n)=1$ 永不收缩，因此
$T_n\not\xrightarrow{\mathbb P}\theta$。例如 $\mathbb P(|T_n-\theta|>1)=2\{1-\Phi(1)\}$ 不随 $n$ 下降。

\section{从分布假设到参数估计}
\label{sec:likelihood-construction}

设模型给出样本的联合概率质量函数或密度 $f_\theta(x_1,\ldots,x_n)$。观测固定后，把它看作参数 $\theta$ 的函数，得到\textbf{似然函数}
\[
 L(\theta;x)=f_\theta(x),\qquad
 \hat\theta_{\rm ML}\in\arg\max_{\theta\in\Theta}L(\theta;x).
\]
概率模型让参数固定、样本变化；似然让已观测样本固定、参数变化。它不是参数的概率分布，通常也不对参数积分为 1。取对数不改变最大点，却把独立样本的乘积变成和，便于求导。最大点可能位于边界，也可能不唯一或不存在，因此求导等于零只是求解的一步。

\subsection{Bernoulli：从十次观测估计成功概率}
若 $X_i$ 独立同分布为 Bernoulli$(p)$，令 $s=\sum_i x_i$，则
\[
 L(p)=p^s(1-p)^{n-s},\qquad
 \ell(p)=s\log p+(n-s)\log(1-p).
\]
当 $0<s<n$，在 $0<p<1$ 内
\[
 \ell'(p)=\frac{s}{p}-\frac{n-s}{1-p}
 =\frac{s-np}{p(1-p)},\qquad
 \ell''(p)=-\frac{s}{p^2}-\frac{n-s}{(1-p)^2}<0.
\]
所以唯一最大点是 $\hat p=s/n=\bar X$。若 $s=0$ 或 $s=n$，分别在允许的边界 $p=0$ 或 $p=1$ 达到最大；只在开区间求导会漏掉这两个样本情形。

对观测 $(1,1,0,1,1,0,1,0,1,1)$，$s=7,n=10$，故 $\hat p=0.7$。似然比较为
\[
 \frac{L(0.7)}{L(0.5)}
 =\frac{0.7^7\,0.3^3}{0.5^{10}}\approx2.277.
\]
它只说这组样本在两个指定参数下的相对支持程度，不是 $p=0.7$ 的后验概率。独立性又给出 $\mathbb E\hat p=p$、$\operatorname{Var}(\hat p)=p(1-p)/n$，由此回到上一节对估计量性质的评价。

\subsection{正态误差为什么导出最小二乘}
给定满列秩 $X\in\mathbb R^{n\times d}$，假设
$y_i=x_i^T\beta+\varepsilon_i$，误差独立服从 $N(0,\sigma^2)$。联合密度为各行正态密度之积，因此
\[
 \ell(\beta,\sigma^2)
 =-\frac n2\log(2\pi)-\frac n2\log\sigma^2
 -\frac{\|y-X\beta\|^2}{2\sigma^2}.
\]
固定任意 $\sigma^2>0$，前两项不随 $\beta$ 变化，最大化最后一项等价于最小化残差平方和。求导得
\[
 \nabla_\beta\ell=\frac{X^T(y-X\beta)}{\sigma^2}=0,
 \qquad \hat\beta=(X^TX)^{-1}X^Ty.
\]
Hessian 为 $-X^TX/\sigma^2\prec0$，所以这个驻点是唯一最大点。再固定 $\hat\beta$，记 $v=\sigma^2$、$S=\|y-X\hat\beta\|^2>0$，则
\[
 \frac{\partial\ell}{\partial v}=-\frac n{2v}+\frac S{2v^2},
 \qquad \hat\sigma^2_{\rm ML}=S/n.
\]
导数在 $v<S/n$ 时为正、在 $v>S/n$ 时为负，故确为最大值。它与无偏估计 $S/(n-d)$ 不同：似然最大化并不以无偏为目标。若 $S=0$，似然随 $v\downarrow0$ 无界，在 $v>0$ 内没有最大点。

沿用第 6 章的三点数据
\[
 X=\begin{pmatrix}1&0\\1&1\\1&2\end{pmatrix},\quad y=(1,2,2)^T,
 \quad X^TX=\begin{pmatrix}3&3\\3&5\end{pmatrix},\quad X^Ty=(5,6)^T.
\]
两条方程相减得 $2\hat\beta_1=1$，代回得 $\hat\beta_0=7/6$。残差为 $(-1/6,1/3,-1/6)^T$，$S=1/6$，所以
\[
 \hat\sigma^2_{\rm ML}=1/18,\qquad s^2=S/(3-2)=1/6.
\]
到这里，概率模型、平方损失和参数估计连在一起。没有正态假设时仍可定义 OLS，也可在外生等条件下研究它的性质，但不能继续把它称作任意误差分布下的极大似然估计。

\subsection{从独立样本到条件似然}
时间序列的联合密度按条件概率分解为
\[
 f_\theta(y_1,\ldots,y_T)=f_\theta(y_1)
 \prod_{t=2}^T f_\theta(y_t\mid y_1,\ldots,y_{t-1}).
\]
独立性只是每个条件密度不再依赖过去的特例。对正态 AR(1)，$y_t=c+\phi y_{t-1}+\eta_t$、$\eta_t\sim N(0,v)$ 独立，给定 $y_1$ 的条件对数似然为
\[
 \ell_c(c,\phi,v)
 =-\frac{T-1}{2}\log(2\pi v)
 -\frac1{2v}\sum_{t=2}^T(y_t-c-\phi y_{t-1})^2.
\]
不加参数约束时，$c,\phi$ 的条件估计就是用滞后值回归当前值；若要求 $|\phi|<1$，还需在该约束内优化。精确似然另含初始密度；平稳 AR(1) 的初始均值为 $c/(1-\phi)$、方差为 $v/(1-\phi^2)$，它们也依赖待估参数。

第 12 章将沿此分解处理变化的条件方差与隐藏状态：GARCH 使每期条件方差不同，Kalman 滤波产生每期创新及其方差。若条件分布未必正态，却使用正态形式作为目标，则称高斯准极大似然；其一致性和标准误依赖额外条件，不能从正态似然公式自动推出。

\section{置信区间怎样由重复抽样产生}

考虑 $X_i\sim N(\theta,\sigma^2)$ 独立，且 $\sigma$ 已知。标准化均值 $Z=\sqrt n(\bar X-\theta)/\sigma$ 精确服从标准正态，故
\[
 \mathbb P_\theta\left(\bar X-1.96\frac\sigma{\sqrt n}
 \leq\theta\leq\bar X+1.96\frac\sigma{\sqrt n}\right)\approx0.95.
\]
这个概率对固定 $\theta$ 和随机端点取值。取 $n=25,\sigma=2,\bar X=1.2$，区间为 $[0.416,1.984]$；检验 $\theta=0$ 时统计量为 3，区间不含零。但如果只挑选 20 次实验中最大的均值，这个枢轴量已被选择改变，不能原样声称选择后的区间有 95\% 覆盖率。

notebook 将生成许多独立数据集，每个数据集画一个区间；先固定报告第一项，再改成报告最大项，比较覆盖率。它还逐步构造 OLS 的残差和稳健协方差，读者可分别改变噪声方差与解释变量的相关性，观察“标准误变化”和“估计目标变化”这两个问题。

\section{渐近近似、检验与经济效应}

一致性要求 $\hat\theta_n\xrightarrow{\mathbb P}\theta$；渐近正态通常写成
\[
 \sqrt n(\hat\theta_n-\theta)\xrightarrow d\mathcal N(0,V).
\]
前者说明误差最终集中，后者给出局部缩放后的近似形状与收敛速率。若两个一致估计量有相同目标和正则条件，才可用较小的渐近方差谈相对效率；不同目标、不同损失或不同模型类不能只比一个标准误。
估计标准误简记为 $\operatorname{se}(\hat\theta)$，它必须对应实际抽样与依赖结构。

\begin{proposition}[渐近正态的使用边界]
若 $Y_i$ 独立同分布、$\mathbb E Y_i=\theta$ 且
$0<\operatorname{Var}(Y_i)=\sigma^2<\infty$，则
\[
 \sqrt n(\bar Y_n-\theta)\xrightarrow d\mathcal N(0,\sigma^2).
\]
若用一致估计 $\widehat\sigma_n^2\to_p\sigma^2$ 替代 $\sigma^2$，Slutsky 定理才允许把
\[
 \frac{\sqrt n(\bar Y_n-\theta)}{\widehat\sigma_n}
\]
近似为标准正态。时间相关、重尾或随机停止会破坏其中的条件；“样本量很大”本身不是标准误一致性的证明。
\end{proposition}

若 $g$ 在 $\theta$ 可微，由一阶 Taylor 展开
\[
 g(\hat\theta_n)-g(\theta)
 =g'(\theta)(\hat\theta_n-\theta)+o_{\mathbb P}(n^{-1/2}),
\]
再用 Slutsky 定理得到 Delta 方法：
\[
 \sqrt n\{g(\hat\theta_n)-g(\theta)\}
 \xrightarrow d\mathcal N\!\left(0,[g'(\theta)]^2V\right).
\]
以 Bernoulli 概率 $p=0.4$、$g(p)=p^2$、$n=250$ 为例，
$\mathbb E(\hat p^2)-p^2=p(1-p)/n=0.000960$；Delta 近似方差为
$4p^2p(1-p)/n=0.0006144$。这说明 plug-in 一致不等于有限样本无偏，也说明导数在边界为零时一阶 Delta 方法可能退化。


一个统计检验由以下要素描述： $H_0,H_1$、统计量、拒绝域、显著性水平和停止规则。检验的\textbf{大小}是零假设参数集合上误拒概率的上确界；在具体备择 $\theta$ 下，功效为
$\mathbb P_\theta(\text{拒绝 }H_0)$。$p$ 值是在 $H_0$ 下取得“至少同样极端”统计量的概率，不是 $H_0$ 为真的概率。

双侧 Wald 检验在
\[
 \left|\frac{\hat\theta-\theta_0}{\widehat{\operatorname{se}}(\hat\theta)}\right|
 >z_{1-\alpha/2}
\]
时拒绝；使用同一标准误时，它与 $(1-\alpha)$ 置信区间不含
$\theta_0$ 对偶。区间是一个随机集合：覆盖率描述重复抽样中覆盖固定真值的比例，不是“参数以 95\% 概率位于已经算出的频率学区间内”。

样本量很大时，极小效应也可显著。效应大小与置信区间帮助判断变化是否具有实际意义；把预测用于交易后，还会遇到样本外衰减、换手、成本、冲击与容量问题。若每看一次结果都用 $0.05$ 阈值，并在十次独立查看中首次显著便停止，则全局零假设下至少一次误拒概率为
$1-0.95^{10}=0.401263$。停止规则属于实验设计，不能在看完路径后删除。

\noindent\textbf{即时练习}\quad
把“$p=0.003$，所以策略有效”改写成有明确条件的结论：注明检验族、停止规则、效应和区间，再列出成本后仍需验证的量。

\section{OLS、识别与不确定性}

对 $y\in\mathbb R^n$、满列秩 $X\in\mathbb R^{n\times k}$，OLS 最小化
$Q(b)=\lVert y-Xb\rVert_2^2$。求梯度得
\[
 \nabla Q(b)=-2X^T(y-Xb).
\]
一阶条件给出正规方程 $X^TX\hat\beta=X^Ty$，所以
\[
 \hat\beta=(X^TX)^{-1}X^Ty,
 \qquad X^T\hat\varepsilon=0.
\]
几何上，$X\hat\beta$ 是 $y$ 到 $\operatorname{col}(X)$ 的正交投影。满秩保证系数表示唯一；数值上仍应优先用 QR/SVD 求解，而不是显式形成逆矩阵。

模型 $y=X\beta+\varepsilon$ 下，给定满列秩设计矩阵时，条件外生
$\mathbb E(\varepsilon\mid X)=0$ 支持有限样本条件无偏。它本身不是完整的
一致性定理；在固定维数的一个常见渐近框架中，还需矩和依赖条件足以应用
大数定律，并要求
\[
 \frac{X^TX}{n}\xrightarrow{p}Q,\qquad Q\succ0,
 \qquad
 \frac{X^T\varepsilon}{n}\xrightarrow{p}0.
\]
于是
$\hat\beta-\beta=(X^TX/n)^{-1}(X^T\varepsilon/n)\xrightarrow{p}0$。
时间序列、面板和高维设计各自还需要相应的依赖、聚类或维数增长条件。
同方差且条件不相关
$\operatorname{Var}(\varepsilon\mid X)=\sigma^2I$ 才给出经典协方差
\[
 \operatorname{Var}(\hat\beta\mid X)=\sigma^2(X^TX)^{-1}.
\]
“系数估计了什么”属于识别问题；“这个系数抽样波动多大”属于协方差估计问题。修正后一项不会自动修正前一项。

前面的正态似然算例已经求出三点回归的系数与残差。下面保持估计公式，改变误差的方差与依赖结构，研究标准误如何变化。

\MFQLead{不同误差结构下的标准误}

令 $x_i^T$ 是第 $i$ 行设计向量。异方差稳健三明治协方差的基本形式是
\[
 (X^TX)^{-1}\left(\sum_i x_ix_i^T\hat\varepsilon_i^2\right)(X^TX)^{-1},
\]
有限样本中还有 HC0--HC3 等杠杆修正。

若误差按时间相关，HAC/Newey--West 用带宽 $L$ 和核权重组合残差得分的滞后协方差；不同带宽会纳入不同范围的相关性，因此可以比较带宽变化时标准误的变化。若观察在公司、日期或事件内相关，簇稳健协方差先在簇内汇总得分，再跨簇求和；有效簇数而非行数控制精度，少簇需要额外修正。若同时存在横截面和时间相关，可能需要双向聚类或更明确的过程模型。

一个解析均值例子把差异量化：$n=100$、单位边际方差下，iid 方差是 $0.010000$；AR(1) 相关系数 $0.6$ 时均值方差为
\[
 \frac{n+2\sum_{h=1}^{n-1}(n-h)0.6^h}{n^2}=0.039250;
\]
簇大小 $5$、簇内相关 $0.4$ 时为
$[1+(5-1)0.4]/100=0.026000$。把它们都当 iid 会系统性低估不确定性。

为看清两种计算，取中心化得分 $(1,1,-1,-1)$。带宽为 $1$ 的 Bartlett--HAC 使用
$\hat\gamma_0+2(1-1/2)\hat\gamma_1$；按标签 $(0,0,1,1)$ 聚类时则先求两个簇得分和。HAC 与簇稳健标准误分别为 $0.559017$ 与 $0.707107$。这里都用除以观察数 4 的协方差约定，未作小样本自由度修正。

\MFQLead{稳健标准误不修复遗漏变量}

真实模型为 $y=\beta_xx+\beta_zz+\varepsilon$，却只回归 $y$ 对 $x$。在相应矩存在且外生时，简单回归斜率的概率极限为
\[
 \operatorname{plim}\hat\beta_x^{\rm short}
 =\beta_x+\beta_z\frac{\operatorname{Cov}(x,z)}{\operatorname{Var}(x)}.
\]
取 $\beta_x=1,\beta_z=2,\operatorname{Cov}(x,z)=0.5$、
$\operatorname{Var}(x)=1$，遗漏后的斜率是 $2$，偏差是 $1$。三明治、HAC 或聚类标准误只会更诚实地描述这个错误目标的抽样波动；它们不会修复遗漏变量、反向因果、未来信息、错误函数形式或选择偏差。

加入控制变量后系数稳定也不是因果证明：控制变量可能测量误差严重、是碰撞点或处于处理之后。因果解释还需要时间顺序、可辩护的识别假设和相应设计。

\section{多重检验与从发现到确认}

同时检验 $m$ 个独立真零假设、每项水平为 $\alpha$ 时，
\[
 \operatorname{FWER}=1-(1-\alpha)^m.
\]
Bonferroni 用阈值 $\alpha/m$，在任意依赖下由并集界控制 FWER。Holm 方法把
第 $j$ 小的值记为 $p_{(j)}$，并将
$p_{(1)}\le\cdots\le p_{(m)}$ 依次与
$\alpha/(m-j+1)$ 比较，遇到首次不拒绝即停止；它同样强控制 FWER，且不比 Bonferroni 更保守。

% Generated by tools/build_figures.py; edit figures/figure-specs.json instead.
\MFQFigure{tex/figures/generated/upper-ch10-multiple-testing.pdf}{同时检验多个假设时，纯噪声也更容易产生极小 p 值；多重检验方法针对整个假设族控制错误。}{upper-ch10-multiple-testing}


Benjamini--Hochberg（BH）控制的是
\[
 \operatorname{FDR}=\mathbb E\!\left[\frac{V}{R\vee1}\right],
\]
其中 $V$ 是错误拒绝数、$R$ 是总拒绝数。找最大 $k$ 使
$p_{(k)}\le kq/m$，拒绝前 $k$ 项。在独立或特定正依赖条件下有标准保证；任意依赖不能机械沿用同一结论。

对固定序列 $(0.001,0.009,0.021,0.040,0.200)$ 与水平 $0.05$，
Bonferroni、Holm、BH 分别拒绝 $(2,2,4)$ 项。BH 多发现两个候选不表示它们可交易，只表示采用了不同错误目标。

在 $m=20,\alpha=0.05$ 且检验独立的全局零假设下，未经校正的 FWER 为 $1-0.95^{20}\approx0.6415$；逐项用 $0.05/20$ 阈值后为 $1-(1-0.05/20)^{20}\approx0.0488$。只发布最显著结果会造成 winner's curse：入选效应向上偏，普通区间也忽略了选择事件。

这个膨胀可在最小例子中精确计算。若两个真实效应均为零，估计噪声独立服从标准正态，则
$\max(Z_1,Z_2)=(Z_1+Z_2+|Z_1-Z_2|)/2$，所以两个独立零效应正态估计中选择较大者，其期望为 $1/\sqrt\pi=0.564190$。增加候选数只会强化选择，不会创造真实 alpha。

\MFQLead{从一个检验结果走到报告结论}

假设某日收益检验得到估计效应 $0.004$、标准误 $0.0015$，未校正双侧 $p$ 值约为 $0.008$。若这是 20 个候选信号中的最佳者，Bonferroni 阈值为 $0.0025$，所以“显著”这一描述在当前族错误率协议下不能直接保留；即使它通过，报告还应给出 $0.004$ 的经济单位、区间、选择前后的样本边界、换手和成本。统计门槛控制的是错误拒绝，不是收益符号、容量或未来稳定性。

\MFQLead{从探索样本到确认样本}

一种分开探索与确认的方法是：先确定假设族、特征与参数网格、停止规则和错误率目标；在探索集上选择；事先确定全部选择；只在未参与选择的时间后移 holdout 上做一次确认。若数据有限，可用嵌套或滚动样本分割，但最外层评估不得回流到选择。样本分割降低估计精度，却换来对选择偏差的隔离。

量化数据还要求时点正确的资产池与特征、重叠收益的 HAC 或分块处理、横截面/日期聚类，以及退市、停牌和制度变化记录。经济解释还涉及成本、换手、冲击、容量与样本外稳定性。FDR 控制不是盈利保证，稳健标准误也不是因果识别。


\noindent\textbf{延伸阅读。}回归与渐近理论参见\cite{hayashi2000econometrics,vanderVaart1998asymptotic}；HAC、横截面回归和数据窥探分别参见\cite{newey1987hac,fama1973risk,white2000reality}。

\section*{章末练习}
\begingroup\small
\begin{enumerate}[leftmargin=*]
  \item \textbf{口述概念：}无偏、一致、渐近正态和渐近有效分别约束什么？给出一个无偏但不一致的估计量。
  \item \textbf{口述概念：}区分检验大小、功效、$p$ 值、置信区间覆盖率与经济效应大小。
  \item \textbf{口述概念：}经典、异方差稳健、HAC 和簇稳健标准误各允许哪类误差结构？它们共同不能修复什么？
  \item \textbf{笔试推导：}证明 MSE 分解，并计算 $0.8\bar X$ 在 $\theta=1$、$\operatorname{Var}(\bar X)=0.25$ 时的偏差、方差和 MSE。
  \item \textbf{笔试推导：}从最小化 $\lVert y-Xb\rVert^2$ 推导正规方程，并证明残差与 $X$ 的每一列正交。
  \item \textbf{笔试推导：}推导简单回归的遗漏变量偏差公式；代入本章固定参数解释为何稳健标准误不能把斜率 $2$ 变回 $1$。

  \item \textbf{数值编程：}重复计算固定第一项与 20 项最大估计的普通正态区间，比较覆盖率。

  \item \textbf{数值编程：}逐步计算 OLS 残差与 HC0 协方差。

  \item \textbf{数值编程：}将样本标准差固定，样本量从 25 增至 100，比较标准误。
  \item \textbf{研究判断：}研究员每天查看一次同一策略的 $p$ 值，显著即停止。设计一个包含登记、停止和确认样本的修复协议。
  \item \textbf{研究判断：}加入行业控制后因子系数稳定且 HC3 显著。列出仍不能声称因果的理由与所需证据。
  \item \textbf{研究判断：}从 500 个候选因子中用 BH 选出 20 个。设计从统计发现到成本后可交易结论的证据链。
\end{enumerate}
\endgroup

\subsection*{基础衔接：独立计算}
已知标准差为 2，样本量从 25 增至 100，95\% 正态区间半宽怎样变化？

\subsection*{分级提示}
\begingroup\small
\begin{enumerate}[leftmargin=*]
  \item 分开有限样本中心、概率收缩、缩放后极限形状与同类估计量的渐近方差。
  \item 大小在零假设下定义，功效在备择下定义；区间覆盖是重复抽样频率。
  \item 先画误差相关矩阵；最后问系数目标本身是否已被正确识别。
  \item 加减 $\mathbb E\hat\theta$，交叉项因中心化消失；收缩会同时改变偏差与方差。
  \item 对 $b$ 求梯度；把一阶条件改写为 $X^T(y-X\hat\beta)=0$。
  \item 将 $z$ 对 $x$ 的线性投影系数代入真实模型；标准误不改变概率极限。
  \item 分别按固定第一项和最大项构造区间，统计覆盖固定真值的比例。
  \item 先求 $e=y-X\hat\beta$，再计算 $X^T\operatorname{diag}(e_i^2)X$ 与两侧逆矩阵。
  \item 标准误与 $1/\sqrt n$ 成正比；分清总体标准差与均值标准误。
  \item 十次独立查看的误拒概率是 $1-0.95^{10}$；确认集不得参与选择。
  \item 检查遗漏变量、碰撞点、处理后控制、未来信息、函数形式与测量误差。
  \item 把检验族、选择、样本外、依赖标准误、效应、成本、容量和限制串成不可回流的流程。
\end{enumerate}
\endgroup


\noindent\textbf{本章要点。}概率模型与损失函数帮助构造估计量，抽样分布描述估计的不确定性。回归系数的目标与标准误的计算需要分别理解。多重检验和模型选择进一步改变推断对象；独立确认样本使选择后的结果能够重新接受检验。